\scanpage{21}
\deleted{$\pi$-torseurs. Pour $j\in J$, le $\pi$-torseur correspondant $P_j$ s'envoie dans $\tau_j=$ fibre de $j\in\tau_\pi$ dans $\tau$ \unsure{puis} la fibre de $P_j$ en $t\in\tau_j$ étant le torseur sur le stabilisateur $\pi_t$ de $t$ dans $\pi$ suivant :
on considère
\diagram{\begin{tikzcd}[ampersand replacement=\&]
\Phi(t):T \arrow[r] \& \pi
\end{tikzcd}}
alors $\pi_t$ centralise}

Considérons la restriction de l'opération de $\mathcal E$ sur $\tau$ à $\pi$, d'où un groupoïde $\langle\tau,\pi\rangle$ dont l'ens des objets est $\tau$, l'ens des $\operatorname{Isom}(t,t')$ étant
\[
\pi_{t\to t'}=\{\lambda\in\pi\mid\lambda t=t'\}.
\]
\deleted{\unsure{Les composantes}} On a
\[
\pi_0\langle\tau,\pi\rangle=\tau/\pi\overset{\mathrm{dfn}}=J
\]
\[
\pi_1(\tau,t)=\pi_t\quad(\text{stabilisateur de }t\text{ dans }\pi)
\]
et on trouve que $\pi_t$ centralise $\Phi(t)(T)\subset\pi$. On suppose \uline{$T$ commutatif} et $\deleted{\pi_t=\Phi_t(T)}\quad\Phi_t:T\xrightarrow{\sim}\pi_t$ --- \unsure{en d'autres termes on suppose donné}, en plus de $\tau$ $\mathcal E$-ens \unsure{donc|de} $\pi$-ens,
\mspage{20}
